\paragraph*{04.04.05.02. Ownership übernehmen und Commitment zeigen}
\kcilabel{para:04.04.05.02_OwnershipUndCommitment}{04.04.05.02}
\label{para:04.04.05.02_OwnershipUndCommitment}

% Task
\par Im komplexen Projektumfeld mit hohen Schnittstellenabhängigkeiten, dynamischer Programmplanung und \gls{pb}-Herausforderungen (\gls{r3}/\gls{s4}) war es entscheidend, Ownership für Gesamtkoordination des \gls{cop} zu übernehmen, inklusive Migrationszeitpunkte. Hohes Engagement förderte Zusammenarbeit und proaktives Handeln.

% Action
\par Entzog Counterpart auf \gls{gu}-Seite Verantwortung für \gls{cop}-Koordination aufgrund unzureichender Planungstiefe und übernahm sie selbst (vgl. \autoref{fig:CutOverPlanGU} und \autoref{fig:R2HErgebnisdokumentationStreamI2C}). Etablierte regelmäßige Abstimmungsmeetings mit Stakeholdern zur Fortschrittsverfolgung und Lösungsentwicklung für \gls{gfa}/\gls{s4} und \gls{bfa}. Nutzte intensive Kommunikation für Akzeptanzförderung.

% Result
\par Ownership-Übernahme und Commitment sicherten \gls{cop}-Koordination und förderten Zusammenarbeit. Resultat: Detailliertere \gls{cop}, abgestimmte Deliverables, schnellere Problemlösung. Produktivsetzung \gls{s4}, Migration \gls{bk} 3800 und \gls{hypc} nahezu reibungsfrei hinsichtlich \gls{ricefw}-Objekte, trotz \gls{pb}-Herausforderungen.

% Reflect
\par Operative \gls{cop}-Übernahme in Detailtiefe nicht mehr möglich als \gls{wsr}-\gls{pm}; künftig Auswirkungsanalyse für Parallelbetrieb und \gls{cop} in allen \gls{wsr}-Umsetzungen.